\chapter{Projective Lightcone and Casimir Equations}
\label{ch:volume-iii-projective-lightcone-casimir}

\section{The Projective Null Cone}

Let \(\R^{D+1,1}\) have coordinates
\[
  X^A=(X^+,X^-,X^\mu),
  \qquad
  \mu=1,\ldots,D,
\]
and bilinear form
\begin{equation}
  X\cdot Y
  =
  -\frac12(X^+Y^-+X^-Y^+)
  +
  \delta_{\mu\nu}X^\mu Y^\nu.
  \label{eq:embedding-bilinear-form}
\end{equation}
The null cone is
\[
  X^2=0.
\]
Physical Euclidean points are rays on this cone:
\[
  X\sim \lambda X,
  \qquad
  \lambda>0.
\]
A convenient section is
\begin{equation}
  X(x)^A=(1,x^2,x^\mu).
  \label{eq:embedding-section}
\end{equation}
Then
\[
  X(x)^2=0,
  \qquad
  X(x)\cdot X(y)=-\frac12(x-y)^2.
  \label{eq:embedding-distance}
\]

\begin{figure}[H]
  \centering
  \resizebox{0.88\textwidth}{!}{%
  \begin{tikzpicture}[x=1cm,y=1cm,every node/.style={font=\scriptsize}]
    \tikzset{
      arrow/.style={-{Latex[length=2mm]},line width=0.8pt},
      insertion/.style={circle,fill=violet!80!black,inner sep=1.8pt}
    }
    \begin{scope}
      \node[font=\small] at (0,2.25) {null cone};
      \draw[thick] (-1.45,0) ellipse (1.25 and 0.34);
      \draw[thick] (-2.70,0)--(-1.45,1.85)--(-0.20,0);
      \draw[thick] (-2.70,0)--(-1.45,-1.85)--(-0.20,0);
      \draw[dashed] (-1.45,0) ellipse (0.55 and 0.15);
      \draw[arrow,blue!65!black] (-1.45,0)--(-0.54,1.35);
      \draw[arrow,blue!65!black] (-1.45,0)--(-0.82,0.94);
      \node[blue!65!black] at (-0.20,1.45) {\(X\)};
      \node[blue!65!black] at (-0.62,0.62) {\(\lambda X\)};
      \node at (-1.45,-2.25) {\(X^2=0,\ X\sim\lambda X\)};
    \end{scope}

    \draw[arrow] (1.20,0)--(2.45,0);
    \node[align=center] at (1.82,0.48) {choose\\section};

    \begin{scope}[shift={(5.20,0)}]
      \node[font=\small] at (0,2.25) {Euclidean section};
      \draw[thick] (-1.80,-1.10)--(1.80,-0.80)--(2.20,1.05)--(-1.50,0.88)--cycle;
      \node[insertion] at (-0.62,-0.25) {};
      \node[insertion] at (0.95,0.42) {};
      \node[violet!80!black] at (-0.85,-0.58) {\(x\)};
      \node[violet!80!black] at (1.20,0.72) {\(y\)};
      \draw[blue!65!black,line width=0.8pt] (-0.62,-0.25)--(0.95,0.42);
      \node[blue!65!black] at (0.02,0.32) {\((x-y)^2\)};
      \node[orange!85!black] at (0.10,1.35)
        {\(-2X(x)\cdot X(y)\)};
    \end{scope}
  \end{tikzpicture}
  }
  \caption{Embedding space realizes a Euclidean point as a ray on the null
  cone.  Inner products of rays reproduce squared distances on a section.}
  \label{fig:projective-lightcone}
\end{figure}

\section{Homogeneous Representatives of Primaries}

A scalar primary of dimension \(\Delta\) is represented in embedding space by
a homogeneous function \(\Phi(X)\) on the null cone:
\begin{equation}
  \Phi(\lambda X)=\lambda^{-\Delta}\Phi(X).
  \label{eq:embedding-scalar-homogeneity}
\end{equation}
The physical operator on the section is
\[
  \mathcal O(x)=\Phi(X(x)).
\]
Conformal transformations are linear \(SO(D+1,1)\) transformations of
\(\R^{D+1,1}\) followed by a return to the chosen section.  The compensating
projective rescaling is exactly the conformal factor in the primary
transformation law.

The scalar two-point function follows from homogeneity and
\eqref{eq:embedding-distance}:
\begin{equation}
  \left\langle\Phi_i(X_1)\Phi_j(X_2)\right\rangle
  =
  \frac{C_{ij}\,\delta_{\Delta_i,\Delta_j}}
       {(-2X_1\cdot X_2)^{\Delta_i}}.
  \label{eq:embedding-two-point}
\end{equation}
Similarly, the scalar three-point function is
\begin{equation}
  \left\langle\Phi_1(X_1)\Phi_2(X_2)\Phi_3(X_3)\right\rangle
  =
  \frac{C_{123}}
  {X_{12}^{(\Delta_1+\Delta_2-\Delta_3)/2}
   X_{23}^{(\Delta_2+\Delta_3-\Delta_1)/2}
   X_{13}^{(\Delta_1+\Delta_3-\Delta_2)/2}},
  \label{eq:embedding-three-point}
\end{equation}
where
\[
  X_{ij}=-2X_i\cdot X_j.
\]
Restricting to the section \(X_i=X(x_i)\) reproduces the flat-space
correlators.

\section{Linear Generators and the Casimir}

The conformal generators act linearly in embedding space:
\begin{equation}
  L_{AB}
  =
  X_A\frac{\partial}{\partial X^B}
  -
  X_B\frac{\partial}{\partial X^A}.
  \label{eq:embedding-generators}
\end{equation}
The quadratic Casimir is
\[
  \mathcal C=-\frac12 L_{AB}L^{AB}.
\]
On a primary of dimension \(\Delta\) and spin \(\ell\), its eigenvalue is
\begin{equation}
  C_{\Delta,\ell}
  =
  \Delta(\Delta-D)+\ell(\ell+D-2).
  \label{eq:conformal-casimir-eigenvalue}
\end{equation}

In a four-point function, the conformal block for exchange of
\((\Delta,\ell)\) in the \(12\) channel is an eigenfunction of the pairwise
Casimir:
\begin{equation}
  \left[
    -\frac12(L_1+L_2)_{AB}(L_1+L_2)^{AB}
    -
    C_{\Delta,\ell}
  \right]
  G_{\Delta,\ell}(u,v)
  =
  0.
  \label{eq:block-casimir-equation-embedding}
\end{equation}
This equation, together with the leading OPE behavior
\eqref{eq:block-leading-radial}, determines the conformal block.

\section{Casimir Equation in Cross-Ratio Coordinates}

For scalar external operators, write the cross ratios as
\[
  u=z\bar z,
  \qquad
  v=(1-z)(1-\bar z).
\]
The Casimir equation takes the form
\begin{equation}
  \left[
    \mathcal D_z+\mathcal D_{\bar z}
    +(D-2)\frac{z\bar z}{z-\bar z}
    \bigl((1-z)\partial_z-(1-\bar z)\partial_{\bar z}\bigr)
  \right]F_{\Delta,\ell}
  =
  \frac12C_{\Delta,\ell}F_{\Delta,\ell}.
  \label{eq:casimir-z-zbar-equation}
\end{equation}
With
\[
  a=-\frac12(\Delta_1-\Delta_2),
  \qquad
  b=\frac12(\Delta_3-\Delta_4),
\]
one convenient convention is
\begin{equation}
  \mathcal D_z
  =
  z^2(1-z)\partial_z^2
  -
  (a+b+1)z^2\partial_z
  -
  ab\,z,
  \label{eq:one-variable-casimir-operator}
\end{equation}
and \(\mathcal D_{\bar z}\) is obtained by replacing
\(z\) with \(\bar z\).  Different normalizations of the four-point
kinematic prefactor conjugate this differential operator by known powers of
\(z,\bar z,1-z,1-\bar z\); the eigenvalue
\eqref{eq:conformal-casimir-eigenvalue} is invariant.

In \(D=2\), the term proportional to \(D-2\) vanishes, and the global block
factorizes into holomorphic and antiholomorphic hypergeometric functions.  In
\(D=4\), the scalar blocks can also be written in closed hypergeometric form.
In odd dimensions, recursive radial methods are usually the most efficient
way to evaluate the blocks to high precision.

\section{Radial Recursion}

Let
\[
  \rho=r\ee^{\ii\alpha}
  =
  \frac{z}{(1+\sqrt{1-z})^2},
  \qquad
  \eta=\cos\alpha.
\]
A radial representation of the block has the form
\begin{equation}
  G_{\Delta,\ell}(u,v)
  =
  (4r)^\Delta h_{\Delta,\ell}(r,\eta).
  \label{eq:block-radial-h}
\end{equation}
The function \(h_{\Delta,\ell}\) is meromorphic in \(\Delta\).  Its poles
occur when a descendant becomes primary, and the residues are fixed by the
corresponding shortening condition.  This gives the radial recursion
\begin{equation}
  h_{\Delta,\ell}
  =
  h_{\infty,\ell}
  +
  \sum_A
  \frac{R_A}{\Delta-\Delta_A^\ast}
  (4r)^{n_A}h_{\Delta_A^\ast+n_A,\ell_A}.
  \label{eq:block-radial-recursion}
\end{equation}
The label \(A\) runs over families of null descendants.  The data
\((\Delta_A^\ast,n_A,\ell_A,R_A)\) are determined by conformal
representation theory.  This recursion makes high-precision bootstrap
calculations possible because \(|\rho|\) is small at the crossing-symmetric
point.
